\IfFileExists{stacks-project.cls}{%
\documentclass{stacks-project}
}{%
\documentclass{amsart}
}

% For dealing with references we use the comment environment
\usepackage{verbatim}
\newenvironment{reference}{\comment}{\endcomment}
%\newenvironment{reference}{}{}
\newenvironment{slogan}{\comment}{\endcomment}
\newenvironment{history}{\comment}{\endcomment}

% For commutative diagrams we use Xy-pic
\usepackage[all]{xy}

% We use 2cell for 2-commutative diagrams.
\xyoption{2cell}
\UseAllTwocells

% We use multicol for the list of chapters between chapters
\usepackage{multicol}

% This is generally recommended for better output
\usepackage{lmodern}
\usepackage[T1]{fontenc}

% For cross-file-references
\usepackage{xr-hyper}

% Package for hypertext links:
\usepackage{hyperref}

% For any local file, say "hello.tex" you want to link to please
% use \externaldocument[hello-]{hello}
\externaldocument[introduction-]{introduction}
\externaldocument[conventions-]{conventions}
\externaldocument[sets-]{sets}
\externaldocument[categories-]{categories}
\externaldocument[topology-]{topology}
\externaldocument[sheaves-]{sheaves}
\externaldocument[sites-]{sites}
\externaldocument[stacks-]{stacks}
\externaldocument[fields-]{fields}
\externaldocument[algebra-]{algebra}
\externaldocument[brauer-]{brauer}
\externaldocument[homology-]{homology}
\externaldocument[derived-]{derived}
\externaldocument[simplicial-]{simplicial}
\externaldocument[more-algebra-]{more-algebra}
\externaldocument[smoothing-]{smoothing}
\externaldocument[modules-]{modules}
\externaldocument[sites-modules-]{sites-modules}
\externaldocument[injectives-]{injectives}
\externaldocument[cohomology-]{cohomology}
\externaldocument[sites-cohomology-]{sites-cohomology}
\externaldocument[dga-]{dga}
\externaldocument[dpa-]{dpa}
\externaldocument[sdga-]{sdga}
\externaldocument[hypercovering-]{hypercovering}
\externaldocument[schemes-]{schemes}
\externaldocument[constructions-]{constructions}
\externaldocument[properties-]{properties}
\externaldocument[morphisms-]{morphisms}
\externaldocument[coherent-]{coherent}
\externaldocument[divisors-]{divisors}
\externaldocument[limits-]{limits}
\externaldocument[varieties-]{varieties}
\externaldocument[topologies-]{topologies}
\externaldocument[descent-]{descent}
\externaldocument[perfect-]{perfect}
\externaldocument[more-morphisms-]{more-morphisms}
\externaldocument[flat-]{flat}
\externaldocument[groupoids-]{groupoids}
\externaldocument[more-groupoids-]{more-groupoids}
\externaldocument[etale-]{etale}
\externaldocument[chow-]{chow}
\externaldocument[intersection-]{intersection}
\externaldocument[pic-]{pic}
\externaldocument[weil-]{weil}
\externaldocument[adequate-]{adequate}
\externaldocument[dualizing-]{dualizing}
\externaldocument[duality-]{duality}
\externaldocument[discriminant-]{discriminant}
\externaldocument[derham-]{derham}
\externaldocument[local-cohomology-]{local-cohomology}
\externaldocument[algebraization-]{algebraization}
\externaldocument[curves-]{curves}
\externaldocument[resolve-]{resolve}
\externaldocument[models-]{models}
\externaldocument[functors-]{functors}
\externaldocument[equiv-]{equiv}
\externaldocument[pione-]{pione}
\externaldocument[etale-cohomology-]{etale-cohomology}
\externaldocument[proetale-]{proetale}
\externaldocument[relative-cycles-]{relative-cycles}
\externaldocument[more-etale-]{more-etale}
\externaldocument[trace-]{trace}
\externaldocument[crystalline-]{crystalline}
\externaldocument[spaces-]{spaces}
\externaldocument[spaces-properties-]{spaces-properties}
\externaldocument[spaces-morphisms-]{spaces-morphisms}
\externaldocument[decent-spaces-]{decent-spaces}
\externaldocument[spaces-cohomology-]{spaces-cohomology}
\externaldocument[spaces-limits-]{spaces-limits}
\externaldocument[spaces-divisors-]{spaces-divisors}
\externaldocument[spaces-over-fields-]{spaces-over-fields}
\externaldocument[spaces-topologies-]{spaces-topologies}
\externaldocument[spaces-descent-]{spaces-descent}
\externaldocument[spaces-perfect-]{spaces-perfect}
\externaldocument[spaces-more-morphisms-]{spaces-more-morphisms}
\externaldocument[spaces-flat-]{spaces-flat}
\externaldocument[spaces-groupoids-]{spaces-groupoids}
\externaldocument[spaces-more-groupoids-]{spaces-more-groupoids}
\externaldocument[bootstrap-]{bootstrap}
\externaldocument[spaces-pushouts-]{spaces-pushouts}
\externaldocument[spaces-chow-]{spaces-chow}
\externaldocument[groupoids-quotients-]{groupoids-quotients}
\externaldocument[spaces-more-cohomology-]{spaces-more-cohomology}
\externaldocument[spaces-simplicial-]{spaces-simplicial}
\externaldocument[spaces-duality-]{spaces-duality}
\externaldocument[formal-spaces-]{formal-spaces}
\externaldocument[restricted-]{restricted}
\externaldocument[spaces-resolve-]{spaces-resolve}
\externaldocument[formal-defos-]{formal-defos}
\externaldocument[defos-]{defos}
\externaldocument[cotangent-]{cotangent}
\externaldocument[examples-defos-]{examples-defos}
\externaldocument[algebraic-]{algebraic}
\externaldocument[examples-stacks-]{examples-stacks}
\externaldocument[stacks-sheaves-]{stacks-sheaves}
\externaldocument[criteria-]{criteria}
\externaldocument[artin-]{artin}
\externaldocument[quot-]{quot}
\externaldocument[stacks-properties-]{stacks-properties}
\externaldocument[stacks-morphisms-]{stacks-morphisms}
\externaldocument[stacks-limits-]{stacks-limits}
\externaldocument[stacks-cohomology-]{stacks-cohomology}
\externaldocument[stacks-perfect-]{stacks-perfect}
\externaldocument[stacks-introduction-]{stacks-introduction}
\externaldocument[stacks-more-morphisms-]{stacks-more-morphisms}
\externaldocument[stacks-geometry-]{stacks-geometry}
\externaldocument[moduli-]{moduli}
\externaldocument[moduli-curves-]{moduli-curves}
\externaldocument[examples-]{examples}
\externaldocument[exercises-]{exercises}
\externaldocument[guide-]{guide}
\externaldocument[desirables-]{desirables}
\externaldocument[coding-]{coding}
\externaldocument[obsolete-]{obsolete}
\externaldocument[fdl-]{fdl}
\externaldocument[index-]{index}

% Theorem environments.
%
\theoremstyle{plain}
\newtheorem{theorem}[subsection]{Theorem}
\newtheorem{proposition}[subsection]{Proposition}
\newtheorem{lemma}[subsection]{Lemma}

\theoremstyle{definition}
\newtheorem{definition}[subsection]{Definition}
\newtheorem{example}[subsection]{Example}
\newtheorem{exercise}[subsection]{Exercise}
\newtheorem{situation}[subsection]{Situation}

\theoremstyle{remark}
\newtheorem{remark}[subsection]{Remark}
\newtheorem{remarks}[subsection]{Remarks}

\numberwithin{equation}{subsection}

% Macros
%
\def\lim{\mathop{\mathrm{lim}}\nolimits}
\def\colim{\mathop{\mathrm{colim}}\nolimits}
\def\Spec{\mathop{\mathrm{Spec}}}
\def\Hom{\mathop{\mathrm{Hom}}\nolimits}
\def\Ext{\mathop{\mathrm{Ext}}\nolimits}
\def\SheafHom{\mathop{\mathcal{H}\!\mathit{om}}\nolimits}
\def\SheafExt{\mathop{\mathcal{E}\!\mathit{xt}}\nolimits}
\def\Sch{\mathit{Sch}}
\def\Mor{\mathop{\mathrm{Mor}}\nolimits}
\def\Ob{\mathop{\mathrm{Ob}}\nolimits}
\def\Sh{\mathop{\mathit{Sh}}\nolimits}
\def\NL{\mathop{N\!L}\nolimits}
\def\CH{\mathop{\mathrm{CH}}\nolimits}
\def\proetale{{pro\text{-}\acute{e}tale}}
\def\etale{{\acute{e}tale}}
\def\QCoh{\mathit{QCoh}}
\def\Ker{\mathop{\mathrm{Ker}}}
\def\Im{\mathop{\mathrm{Im}}}
\def\Coker{\mathop{\mathrm{Coker}}}
\def\Coim{\mathop{\mathrm{Coim}}}

% Boxtimes
%
\DeclareMathSymbol{\boxtimes}{\mathbin}{AMSa}{"02}

%
% Macros for moduli stacks/spaces
%
\def\QCohstack{\mathcal{QC}\!\mathit{oh}}
\def\Cohstack{\mathcal{C}\!\mathit{oh}}
\def\Spacesstack{\mathcal{S}\!\mathit{paces}}
\def\Quotfunctor{\mathrm{Quot}}
\def\Hilbfunctor{\mathrm{Hilb}}
\def\Curvesstack{\mathcal{C}\!\mathit{urves}}
\def\Polarizedstack{\mathcal{P}\!\mathit{olarized}}
\def\Complexesstack{\mathcal{C}\!\mathit{omplexes}}
% \Pic is the operator that assigns to X its picard group, usage \Pic(X)
% \Picardstack_{X/B} denotes the Picard stack of X over B
% \Picardfunctor_{X/B} denotes the Picard functor of X over B
\def\Pic{\mathop{\mathrm{Pic}}\nolimits}
\def\Picardstack{\mathcal{P}\!\mathit{ic}}
\def\Picardfunctor{\mathrm{Pic}}
\def\Deformationcategory{\mathcal{D}\!\mathit{ef}}
\usepackage{cleveref}
\crefname{originallemma}{original lemma}{original lemmas}
\Crefname{originallemma}{Original lemma}{Original lemmas}
\makeatletter
\newlabel{curves-lemma-genus-in-nodal-family-of-curves}{{0GKA}{source}{}{sourcepointer-1}{}}
\newlabel{curves-lemma-genus-in-nodal-family-of-curves@cref}{{[originallemma][0][]0GKA}{[1][0][]source}{}{}{}}
\makeatother
\crefname{originallemma}{original lemma}{original lemmas}
\Crefname{originallemma}{Original lemma}{Original lemmas}
\makeatletter
\newlabel{etale-cohomology-lemma-characterize-locally-constant-module}{{0GKC}{source}{}{sourcepointer-2}{}}
\newlabel{etale-cohomology-lemma-characterize-locally-constant-module@cref}{{[originallemma][0][]0GKC}{[1][0][]source}{}{}{}}
\makeatother
\crefname{originallemma}{original lemma}{original lemmas}
\Crefname{originallemma}{Original lemma}{Original lemmas}
\makeatletter
\newlabel{etale-cohomology-lemma-proper-base-change-stalk}{{0DDF}{source}{}{sourcepointer-3}{}}
\newlabel{etale-cohomology-lemma-proper-base-change-stalk@cref}{{[originallemma][0][]0DDF}{[1][0][]source}{}{}{}}
\makeatother
\crefname{originallemma}{original lemma}{original lemmas}
\Crefname{originallemma}{Original lemma}{Original lemmas}
\makeatletter
\newlabel{etale-cohomology-lemma-sections-upstairs}{{0A3H}{source}{}{sourcepointer-4}{}}
\newlabel{etale-cohomology-lemma-sections-upstairs@cref}{{[originallemma][0][]0A3H}{[1][0][]source}{}{}{}}
\makeatother
\crefname{originallemma}{original lemma}{original lemmas}
\Crefname{originallemma}{Original lemma}{Original lemmas}
\makeatletter
\newlabel{etale-cohomology-lemma-sp-isom-proper-loc-cst-torsion}{{0GKD}{source}{}{sourcepointer-5}{}}
\newlabel{etale-cohomology-lemma-sp-isom-proper-loc-cst-torsion@cref}{{[originallemma][0][]0GKD}{[1][0][]source}{}{}{}}
\makeatother
\crefname{originalreference}{original reference}{original references}
\Crefname{originalreference}{Original reference}{Original references}
\makeatletter
\newlabel{etale-cohomology-section-artin-schreier}{{0A3J}{source}{}{sourcepointer-6}{}}
\newlabel{etale-cohomology-section-artin-schreier@cref}{{[originalreference][0][]0A3J}{[1][0][]source}{}{}{}}
\makeatother
\crefname{originaltheorem}{original theorem}{original theorems}
\Crefname{originaltheorem}{Original theorem}{Original theorems}
\makeatletter
\newlabel{etale-cohomology-theorem-proper-base-change}{{095T}{source}{}{sourcepointer-7}{}}
\newlabel{etale-cohomology-theorem-proper-base-change@cref}{{[originaltheorem][0][]095T}{[1][0][]source}{}{}{}}
\makeatother
\crefname{originaltheorem}{original theorem}{original theorems}
\Crefname{originaltheorem}{Original theorem}{Original theorems}
\makeatletter
\newlabel{etale-cohomology-theorem-vanishing-affine-curves}{{03SC}{source}{}{sourcepointer-8}{}}
\newlabel{etale-cohomology-theorem-vanishing-affine-curves@cref}{{[originaltheorem][0][]03SC}{[1][0][]source}{}{}{}}
\makeatother
\crefname{originallemma}{original lemma}{original lemmas}
\Crefname{originallemma}{Original lemma}{Original lemmas}
\makeatletter
\newlabel{lemma-frobenius-linear-on-vb}{{0GKU}{source}{}{sourcepointer-9}{}}
\newlabel{lemma-frobenius-linear-on-vb@cref}{{[originallemma][0][]0GKU}{[1][0][]source}{}{}{}}
\makeatother
\crefname{originallemma}{original lemma}{original lemmas}
\Crefname{originallemma}{Original lemma}{Original lemmas}
\makeatletter
\newlabel{limits-lemma-descend-dimension-d}{{0EY2}{source}{}{sourcepointer-10}{}}
\newlabel{limits-lemma-descend-dimension-d@cref}{{[originallemma][0][]0EY2}{[1][0][]source}{}{}{}}
\makeatother
\crefname{originallemma}{original lemma}{original lemmas}
\Crefname{originallemma}{Original lemma}{Original lemmas}
\makeatletter
\newlabel{limits-lemma-descend-finite-presentation}{{01ZM}{source}{}{sourcepointer-11}{}}
\newlabel{limits-lemma-descend-finite-presentation@cref}{{[originallemma][0][]01ZM}{[1][0][]source}{}{}{}}
\makeatother
\crefname{originallemma}{original lemma}{original lemmas}
\Crefname{originallemma}{Original lemma}{Original lemmas}
\makeatletter
\newlabel{limits-lemma-descend-smooth}{{0C0C}{source}{}{sourcepointer-12}{}}
\newlabel{limits-lemma-descend-smooth@cref}{{[originallemma][0][]0C0C}{[1][0][]source}{}{}{}}
\makeatother
\crefname{originallemma}{original lemma}{original lemmas}
\Crefname{originallemma}{Original lemma}{Original lemmas}
\makeatletter
\newlabel{limits-lemma-eventually-proper}{{081F}{source}{}{sourcepointer-13}{}}
\newlabel{limits-lemma-eventually-proper@cref}{{[originallemma][0][]081F}{[1][0][]source}{}{}{}}
\makeatother
\crefname{originallemma}{original lemma}{original lemmas}
\Crefname{originallemma}{Original lemma}{Original lemmas}
\makeatletter
\newlabel{more-morphisms-lemma-proper-flat-geom-red}{{0E0N}{source}{}{sourcepointer-14}{}}
\newlabel{more-morphisms-lemma-proper-flat-geom-red@cref}{{[originallemma][0][]0E0N}{[1][0][]source}{}{}{}}
\makeatother
\title{Constructibility of R\textasciicircum{}qf\_* of the constant mod-l sheaf for a proper smooth family of curves with geometrically connected fibres}
\author{The Stacks Project Authors}
\date{stacks-project:more-etale.tex:local-snapshot-20261012; GNU FDL 1.2 or later}
\begin{document}
\maketitle
\setcounter{subsection}{1}
\noindent\textit{Editorial note.} All delivered mathematical text is the author's own, copied from the pinned Stacks Project chapter file more-etale.tex; no statement or proof step was rewritten, reordered or removed. The counted span is the author environment and its own proof environment. References inside the span are handed to the assembler as navigation stubs that carry the canonical tag of the cited tag only; whether each of them is load-bearing is recorded on the stub itself.\par\medskip
\noindent\textit{Allowed standard prerequisites (editorial).}\begin{itemize}
\item Proper base change theorem (Etale Cohomology, Theorem etale-cohomology-theorem-proper-base-change, canonical tag 095T). cited at proof lines 4526-4530. Base change compatibility of R\textasciicircum{}qf\_* underwrites the split of S into the l invertible and l zero parts, and its stalk form drives both cases.
\item Artin-Schreyer sequence (Etale Cohomology, Section etale-cohomology-section-artin-schreier, canonical tag 0A3J). cited at proof lines 4560-4566. Supplies the sequence 0 -> Z/pZ -> O\_X -> O\_X -> 0 from which R\textasciicircum{}1f\_* of the constant sheaf is computed in case (b).
\item Constructibility of Frobenius kernels and cokernels on vector bundles (Lemma lemma-frobenius-linear-on-vb, canonical tag 0GKU). cited at proof line 4579. The terminal step of case (b): the short exact sequence exhibits R\textasciicircum{}1f\_* as built from Coker(F - 1) and Ker(F - 1) on locally free sheaves, to which the lemma applies.
\end{itemize}\medskip
\ifdefined\hypertarget\hypertarget{sourcepointer-1}{}\fi
\noindent\textit{Source reference.} Original lemma 0GKA: sources/stacks-chow/tags.txt, source lines 20680--20680. Cited at proof lines 4567-4570: f\_*O\_X = O\_S and R\textasciicircum{}1f\_*O\_X is finite locally free of rank the genus of the fibres of X -> S.\par\medskip
\ifdefined\hypertarget\hypertarget{sourcepointer-2}{}\fi
\noindent\textit{Source reference.} Original lemma 0GKC: sources/stacks-chow/tags.txt, source lines 20682--20682. Cited at proof lines 4547-4549: finite sheaves whose specialization maps are isomorphisms are finite locally constant, finishing case (a).\par\medskip
\ifdefined\hypertarget\hypertarget{sourcepointer-3}{}\fi
\noindent\textit{Source reference.} Original lemma 0DDF: sources/stacks-chow/tags.txt, source lines 16765--16765. Cited at proof lines 4537-4543: identifies the stalks of R\textasciicircum{}qf\_*Z/lZ with fibre cohomology in cases (a) and (b).\par\medskip
\ifdefined\hypertarget\hypertarget{sourcepointer-4}{}\fi
\noindent\textit{Source reference.} Original lemma 0A3H: sources/stacks-chow/tags.txt, source lines 12740--12740. Cited at proof lines 4556-4558: f\_*Z/pZ = Z/pZ in case (b) since the fibres are geometrically connected.\par\medskip
\ifdefined\hypertarget\hypertarget{sourcepointer-5}{}\fi
\noindent\textit{Source reference.} Original lemma 0GKD: sources/stacks-chow/tags.txt, source lines 20683--20683. Cited at proof lines 4544-4546: all specialization maps are isomorphisms in case (a).\par\medskip
\ifdefined\hypertarget\hypertarget{sourcepointer-6}{}\fi
\noindent\textit{Source reference.} Original reference 0A3J: sources/stacks-chow/tags.txt, source lines 12742--12742. Cited at proof line 4566: source of the Artin-Schreyer sequence 0 -> Z/pZ -> O\_X -> O\_X -> 0 used in case (b).\par\medskip
\ifdefined\hypertarget\hypertarget{sourcepointer-7}{}\fi
\noindent\textit{Source reference.} Original theorem 095T: sources/stacks-chow/tags.txt, source lines 11596--11596. Cited at proof lines 4526-4530: formation of R\textasciicircum{}qf\_*Z/lZ commutes with arbitrary base change, enabling the split into the two cases.\par\medskip
\ifdefined\hypertarget\hypertarget{sourcepointer-8}{}\fi
\noindent\textit{Source reference.} Original theorem 03SC: sources/stacks-chow/tags.txt, source lines 4963--4963. Cited at proof lines 4540-4543: finiteness of the stalks of R\textasciicircum{}qf\_*Z/lZ in case (a), and by the same references in case (b).\par\medskip
\ifdefined\hypertarget\hypertarget{sourcepointer-9}{}\fi
\noindent\textit{Source reference.} Original lemma 0GKU: sources/stacks-chow/tags.txt, source lines 20699--20699. Cited at proof line 4579: concludes case (b) by applying the lemma to the Coker and Ker of Frobenius-minus-one appearing in the short exact sequence.\par\medskip
\ifdefined\hypertarget\hypertarget{sourcepointer-10}{}\fi
\noindent\textit{Source reference.} Original lemma 0EY2: sources/stacks-chow/tags.txt, source lines 18677--18677. Cited at proof lines 4510-4515: after increasing i, f\_i has relative dimension 1.\par\medskip
\ifdefined\hypertarget\hypertarget{sourcepointer-11}{}\fi
\noindent\textit{Source reference.} Original lemma 01ZM: sources/stacks-chow/tags.txt, source lines 2599--2599. Cited at proof lines 4505-4509: descend f to a finitely presented morphism f\_i : X\_i -> S\_i over a finite type Z-subalgebra.\par\medskip
\ifdefined\hypertarget\hypertarget{sourcepointer-12}{}\fi
\noindent\textit{Source reference.} Original lemma 0C0C: sources/stacks-chow/tags.txt, source lines 15082--15082. Cited at proof lines 4510-4515: after increasing i, f\_i is smooth.\par\medskip
\ifdefined\hypertarget\hypertarget{sourcepointer-13}{}\fi
\noindent\textit{Source reference.} Original lemma 081F: sources/stacks-chow/tags.txt, source lines 10215--10215. Cited at proof lines 4510-4515: after increasing i, f\_i is proper.\par\medskip
\ifdefined\hypertarget\hypertarget{sourcepointer-14}{}\fi
\noindent\textit{Source reference.} Original lemma 0E0N: sources/stacks-chow/tags.txt, source lines 17543--17543. Cited at proof lines 4516-4518: the open U\_i subset S\_i over which the fibres of f\_i are geometrically connected, used to enter the Noetherian case with the hypothesis in place.\par\medskip
\begin{lemma}
\label{lemma-proper-smooth-family-curves}
Let $f : X \to S$ be a proper smooth morphism of schemes with
geometrically connected fibres of dimension $1$. Let $\ell$
be a prime number. Then $R^qf_*\underline{\mathbf{Z}/\ell\mathbf{Z}}$
is a constructible.
\end{lemma}

\begin{proof}
We may assume $S$ is affine. Say $S = \Spec(A)$. Then, if we write
$A = \bigcup A_i$ as the union of its finite type $\mathbf{Z}$-subalgebras,
we can find an $i$ and a morphism $f_i : X_i \to S_i = \Spec(A_i)$ of finite
type whose base change to $S$ is $f : X \to S$, see
Limits, Lemma \ref{limits-lemma-descend-finite-presentation}.
After increasing $i$ we may assume $f_i : X_i \to S_i$ is
smooth, proper, and of relative dimension $1$, see
Limits, Lemmas
\ref{limits-lemma-eventually-proper}
\ref{limits-lemma-descend-smooth}, and
\ref{limits-lemma-descend-dimension-d}.
By More on Morphisms, Lemma \ref{more-morphisms-lemma-proper-flat-geom-red}
we obtain an open subscheme $U_i \subset S_i$ such that the fibres
of $f_i : X_i \to S_i$ over $U_i$ are geometrically connected.
Then $S \to S_i$ maps into $U_i$. We may replace $X \to S$
by $f_i : f_i^{-1}(U_i) \to U_i$ to reduce to the case
discussed in the next paragraph.

\medskip\noindent
Assume $S$ is Noetherian. We may write $S = U \cup Z$ where $U$ is the
open subscheme defined by the nonvanishing of $\ell$ and
$Z = V(\ell) \subset S$. Since the formation of
$R^qf_*\underline{\mathbf{Z}/\ell\mathbf{Z}}$ commutes with arbitrary
base change (\'Etale Cohomology, Theorem
\ref{etale-cohomology-theorem-proper-base-change}),
it suffices to prove the result over $U$ and over $Z$.
Thus we reduce to the following two cases: (a) $\ell$
is invertible on $S$ and (b) $\ell$ is zero on $S$.

\medskip\noindent
Case (a). We claim that in this case the sheaves
$R^qf_*\underline{\mathbf{Z}/\ell\mathbf{Z}}$ are
finite locally constant on $S$. First, by proper base change
(in the form of \'Etale Cohomology, Lemma
\ref{etale-cohomology-lemma-proper-base-change-stalk})
and by finiteness 
(\'Etale Cohomology, Theorem
\ref{etale-cohomology-theorem-vanishing-affine-curves})
we see that the stalks of $R^qf_*\underline{\mathbf{Z}/\ell\mathbf{Z}}$
are finite. By \'Etale Cohomology, Lemma
\ref{etale-cohomology-lemma-sp-isom-proper-loc-cst-torsion}
all specialization maps are isomorphisms.
We conclude the claim holds by
\'Etale Cohomology, Lemma
\ref{etale-cohomology-lemma-characterize-locally-constant-module}.

\medskip\noindent
Case (b). Here $\ell = p$ is a prime and $S$ is a scheme over
$\Spec(\mathbf{F}_p)$. By the same references as above we already
know that the stalks of $R^qf_*\underline{\mathbf{Z}/p\mathbf{Z}}$
are finite and zero for $q \geq 2$. It follows from
\'Etale Cohomology, Lemma \ref{etale-cohomology-lemma-sections-upstairs}
that $f_*\underline{\mathbf{Z}/p\mathbf{Z}} =
\underline{\mathbf{Z}/p\mathbf{Z}}$. It remains to prove
that $R^1f_*\underline{\mathbf{Z}/p\mathbf{Z}}$ is constructible.
Consider the Artin-Schreyer sequence
$$
0 \to \underline{\mathbf{Z}/p\mathbf{Z}} \to \mathcal{O}_X
\xrightarrow{F - 1} \mathcal{O}_X
\to 0
$$
See \'Etale Cohomology, Section \ref{etale-cohomology-section-artin-schreier}.
Recall that $f_*\mathcal{O}_X = \mathcal{O}_S$ and
$R^1f_*\mathcal{O}_X$ is a finite locally free $\mathcal{O}_S$-module
of rank equal to the genera of the fibres of $X \to S$, see
Algebraic Curves, Lemma \ref{curves-lemma-genus-in-nodal-family-of-curves}.
We conclude that we have a short exact sequence
$$
0 \to
\Coker(F - 1 : \mathcal{O}_S \to \mathcal{O}_S)
\to R^1f_*\underline{\mathbf{Z}/p\mathbf{Z}} \to
\Ker(F - 1 : R^1f_*\mathcal{O}_X \to R^1f_*\mathcal{O}_X)
\to 0
$$
Applying Lemma \ref{lemma-frobenius-linear-on-vb} we win.
\end{proof}

\end{document}
